\documentclass[border=10pt]{standalone}
\usepackage{fontspec}
\AtBeginDocument{\hyphenation{zenodo}}% never hyphenate "zenodo" inside a DOI
\setmainfont{Noto Serif}[Ligatures=TeXOff]
\XeTeXgenerateactualtext=1
\newfontfamily\bengalifont{Noto Serif Bengali}[Script=Bengali, NFSSFamily=bnfam]
\usepackage[Latin,Bengali]{ucharclasses}
% group-free font switching (as in header-common.tex): robust inside braces, cells and quotes
\makeatletter
\def\BN@saved{\rmdefault}
\def\BN@bn{bnfam}
\newcommand\BNin{\edef\BN@tmp{\f@family}\ifx\BN@tmp\BN@bn\else\let\BN@saved\BN@tmp\fi\fontfamily{bnfam}\selectfont}
\newcommand\BNout{\fontfamily{\BN@saved}\selectfont}
\makeatother
\setTransitionsFor{Bengali}{\BNin}{\BNout}
% Typographic quotes and dashes are not in a ucharclasses class (class 0), so a Bangla word right
% after “ or ‘ would not trigger the Bangla font. Treat them as Latin so that every switch fires.
\XeTeXcharclass"2018=\XeTeXcharclass`A
\XeTeXcharclass"2019=\XeTeXcharclass`A
\XeTeXcharclass"201C=\XeTeXcharclass`A
\XeTeXcharclass"201D=\XeTeXcharclass`A
\XeTeXcharclass"2013=\XeTeXcharclass`A
\XeTeXcharclass"2014=\XeTeXcharclass`A
\XeTeXcharclass"2026=\XeTeXcharclass`A
\usepackage{newunicodechar}
\newunicodechar{।}{{\bengalifont\char"0964}}
\newunicodechar{॥}{{\bengalifont\char"0965}}

\usepackage[table]{xcolor}
\usepackage{booktabs,tabularx,array}
\definecolor{bapilrule}{HTML}{1F2937}
\definecolor{bapilstripe}{HTML}{F3F4F6}
\renewcommand{\arraystretch}{1.28}
\setlength{\tabcolsep}{5pt}

\begin{document}
\begin{minipage}{17cm}
{\large\bfseries সারণি ২৫।}\ {\large ডেটা-মানের মাত্রা, মাপকাঠি ও লক্ষ্য}\par\vspace{7pt}
\small
\rowcolors{2}{bapilstripe}{white}
\begin{tabularx}{\linewidth}{@{}>{\hsize=0.652\hsize\raggedright\arraybackslash}X>{\hsize=1.478\hsize\raggedright\arraybackslash}X>{\hsize=0.870\hsize\raggedright\arraybackslash}X@{}}
\arrayrulecolor{bapilrule}\toprule
\rowcolor{white}\textbf{মাত্রা} & \textbf{মাপকাঠি} & \textbf{পাইলট লক্ষ্য} \\
\midrule
উৎস-তথ্য & পূর্ণ উৎস, তারিখ, চেকসাম ও পুনর্ব্যবহারের শর্তযুক্ত নথির অনুপাত & ১০০\% \\
টেক্সটের নির্ভুলতা & হাতে-যাচাইকৃত নমুনায় OCR-এর অক্ষর-ত্রুটির হার & সূচিবদ্ধ পৃষ্ঠায় ৩\%-এর নিচে \\
সম্পূর্ণতা & নির্ধারিত সংগ্রহের গৃহীত অংশ (যেমন বলবৎ সব আইন) & লক্ষ্য-সংগ্রহের অন্তত ৯০\% \\
সময়োপযোগিতা & প্রকাশ থেকে অনুসন্ধানে প্রাপ্যতা পর্যন্ত দিনের মধ্যমা & পর্যবেক্ষিত উৎসে ৭ দিনের কম \\
সামঞ্জস্য & মেটাডেটা স্কিমায় বৈধ রেকর্ডের অনুপাত & ১০০\% \\
অনন্যতা & শনাক্ত ও একীভূত দ্বৈত নথির অনুপাত & দ্বৈততা ১\%-এর নিচে \\
ভাষাগত সমতা & যেখানে উভয় ভাষা আছে সেখানে বাংলা ও ইংরেজি শিরোনামযুক্ত নথির অনুপাত & অন্তত ৯৫\% \\
\bottomrule
\end{tabularx}
\par\vspace{5pt}{\footnotesize\itshape পাইলট-লক্ষ্য। প্রতিটি কর্পাস রিলিজে কমপক্ষে ৫০ পৃষ্ঠার স্তরভিত্তিক নমুনায় দুজন যাচাইকারীর মাধ্যমে OCR-এর অক্ষর-ত্রুটির হার পরিমাপ করা হয়।}
\par\vspace{3pt}{\scriptsize\color{gray} উৎস: BAPIL টেকনিক্যাল ব্লুপ্রিন্ট, সংস্করণ ১.০ (২০২৬)। CC BY 4.0।}
\end{minipage}
\end{document}
